\documentclass[../principal.tex]{subfiles}
\graphicspath{{\subfix{../imagenes/}}}

\begin{document}

\chapter{Placa de pruebas}

La placa de pruebas recibe varios nombres según país.

En Estados Unidos se conoce como \textit{breadboard}, en Chile recibe el nombre de  \textit{protoboard}.

Las columnas van de arriba a abajo, y las filas van de izquierda a derecha.

\section{Placa de pruebas de 400 puntos}

Un tamaño fácil de encontrar es el de 400 puntos donde se disponen 10 columnas (abcde y fghij) y 40 filas, de 1 a 40, para un total de 400 puntos donde hacer conexiones.

Por su construcción, las columnas no tienen conexión entre ellas, son independientes entre sí, mientras que las filas están conectadas por abajo por un cable.

La misma fila son dos lugares eléctricos separados por un surco, entonces en la fila 5, tenemos dos lugares eléctricos:

\begin{itemize}
    \item 5a
    \item 5b
    \item 5c
    \item 5d
    \item 5e
\end{itemize}

\begin{itemize}
    \item 5f
    \item 5g
    \item 5h
    \item 5i
    \item 5j
\end{itemize}

\section{Circuitos integrados en placas de pruebas}

Los circuitos integrados en formato DIP se usan en paralelo al surco central.

\newpage

\end{document}

